\section{Fazit}

In diesem Versuch wurden anhand des Frank-Hertz-Versuchs und der Energieausfspaltung aufgrund des Zeeman-Effekts von Atomen im homogenen Magnetfeld die Quantisierung der Energiezustände von Atomen untersucht.

Im ersten Versuchsteil wurde die Zeeman-Aufspaltung an der roten $\lambda_0=\SI{643.847}{\nano\meter}$ von Cadmium untersucht. Die erwarteten Aufspaltungen in die $\sigma_\pm$ und $\pi$ Komponenten konnten in transversaler und logitudinaler Konfiguration wie erwartet beobachtet werden. Aus der Winkelaufspaltung wurde mithilfe von Anpassung von Gaußfunktionen die Energieaufspaltung der Linienkomponenten die Größe des Bohrschen Magnetons auf $\mu_B=\SI{6.140(33)e-5}{\eV\per\tesla}$ bestimmt. Das ist eine Abweichung von $6.1\%$ vom Literaturwert und damit noch eine gute Annäherung.

Im zweiten Versuchsteil wurde der Franck-Hertz-Versuch an Quecksilber durchgeführt. Als Anregungsenergie aus den Abständen der Resonanzkurvenspitzen wurde $\Delta E = \SI{4.816(20)}{\electronvolt}$ bestimmt, was der Erwartung des dominierenden Übergangs $\ce{^1S_0} \rightarrow \ce{^3P_1}$ und $\ce{^1S_0} \rightarrow \ce{^3P_0}$ entspricht. Bei steigender Temperatur zeigte sich wie erwartet ein Abfallen der Anodenstromamplitude. Bei steigender Gegenspannung zeigt sich ein stärkerer Abfall des Stroms zwischen den Maxima, da mehr Elektronen nach einem elastischen Stoß die Anode nicht mehr erreichen können. Dies entspricht ebenfalls den Erwartungen.
